% answers_short.tex —— 1999 年·全国卷·理科 · 简短答案速查
% 由 tools/make_short_answers.py 自 answers.tex 的 \gksolhead/\gkanswer 提取，
% 并与 1999/1999_数学真题_全国理_答案要点.md 逐题比对后生成，勿手改（改源 MD 后重跑）。

\begin{gknotice}
\textbf{使用说明}：本节只给最终答案便于快速核对，完整推导见解析版。答案为编写者独立推导，\textbf{非官方标准答案}。
\end{gknotice}

\gkpart{题目}
\begin{itemize}[leftmargin=2.6em]
  \item[\textbf{1.}] C
  \item[\textbf{2.}] A
  \item[\textbf{3.}] A
  \item[\textbf{4.}] C
  \item[\textbf{5.}] B
  \item[\textbf{6.}] B
  \item[\textbf{7.}] B
  \item[\textbf{8.}] A
  \item[\textbf{9.}] C
  \item[\textbf{10.}] D
  \item[\textbf{11.}] B
  \item[\textbf{12.}] D
  \item[\textbf{13.}] D
  \item[\textbf{14.}] C
  \item[\textbf{15.}] $\dfrac12$
  \item[\textbf{16.}] 6
  \item[\textbf{17.}] $[9,+\infty)$
  \item[\textbf{18.}] 以①③④为条件，②为结论。
  \item[\textbf{19.}] 令 $u=\log_a x$，则 $u\in\left(\dfrac12,\dfrac34\right)\cup(1,+\infty)$；当 $a>1$ 时，$x\in\left(a^{1/2},a^{3/4}\right)\cup(a,+\infty)$；当 $0<a<1$ 时，$x\in(0,a)\cup\left(a^{3/4},a^{1/2}\right)$。
  \item[\textbf{20.}] 最大值为 $\arctan\dfrac{\sqrt6}{12}$，对应 $\theta=\arctan\dfrac{\sqrt6}{2}$。
  \item[\textbf{21.}] (1) $\dfrac{a^2}{\sqrt2}$；(2) $\sqrt2\,a$；(3) $\dfrac{a^3}{2\sqrt2}$
  \item[\textbf{22.}] (1) $\left\lceil\log_{1-r_0}\dfrac{\beta}{\alpha}\right\rceil$ 对轧辊（$0<\beta<\alpha$，$0<r_0<1$）；(2) $L_1=3125\ \mathrm{mm}$，$L_2=2500\ \mathrm{mm}$，$L_3=2000\ \mathrm{mm}$
  \item[\textbf{23.}] (1) $x_1=1$，$x_2=1+\dfrac1b$，$x_n=\displaystyle\sum_{k=0}^{n-1}b^{-k}=\dfrac b{b-1}\left(1-b^{-n}\right)$；(2) $f(x)=n+b^n(x-x_n)$（$x_n\le x\le x_{n+1}$，$n=0,1,2,\ldots$）；$b>1$ 时定义域为 $\left[0,\dfrac b{b-1}\right)$，$0<b<1$ 时定义域为 $[0,+\infty)$；(3) $b>1$ 时 $x>1$ 有 $f(x)>x$，$0<b<1$ 时 $x>1$ 有 $f(x)<x$，所以没有横坐标大于 1 的交点
  \item[\textbf{24.}] 轨迹方程为 $(a-1)x^2+2ax-(a+1)y^2=0$；$a>1$ 时为双曲线，$a=1$ 时为抛物线 $x=y^2$，$0<a<1$ 时为椭圆。
  \item[\textbf{25.}] (1) $b_n=3n-2$；(2) $a>1$ 时 $S_n>\dfrac13\log_a b_{n+1}$；$0<a<1$ 时 $S_n<\dfrac13\log_a b_{n+1}$
\end{itemize}
